\documentclass[12pt]{article}
\usepackage[utf8]{inputenc}
\usepackage[T1]{fontenc}
\usepackage[margin=1in]{geometry}
\usepackage{setspace}
\usepackage{xcolor}
\usepackage[colorlinks=true, urlcolor=blue!60!black, linkcolor=blue!60!black]{hyperref}
\setlength{\parindent}{2em}
\setlength{\parskip}{1.25em}
\renewcommand{\baselinestretch}{1.0}
\title{\textit{Tennessee Wine \& Spirits Retailers Ass'n v. Thomas}, 588 U.S. 504 (2019)}
\author{Case Brief}
\date{}
\begin{document}
\maketitle

\section{Case Metadata}

\begin{itemize}
    \item \textbf{Case Name:} \textit{Tennessee Wine \& Spirits Retailers Ass'n v. Thomas}, 588 U.S. 504 (2019)
    \item \textbf{Docket No.:} 18-96
    \item \textbf{Dates:} Argued January 16, 2019; Decided June 26, 2019
    \item \textbf{Lower Court:} Certiorari to the U.S. Court of Appeals for the Sixth Circuit (\textit{Byrd v. Tennessee Wine \& Spirits Retailers Ass'n}, 883 F.3d 608 (6th Cir. 2018), affirming 259 F. Supp. 3d 785 (M.D. Tenn. 2017))
    \item \textbf{Vote:} 7--2
    \item \textbf{Majority Opinion:} Justice Samuel A. Alito, Jr., joined by Chief Justice John G. Roberts, Jr. and Justices Ruth Bader Ginsburg, Stephen G. Breyer, Sonia Sotomayor, Elena Kagan, and Brett M. Kavanaugh
    \item \textbf{Dissenting Opinion:} Justice Neil M. Gorsuch, joined by Justice Clarence Thomas
    \item \textbf{Constitutional and Statutory Provisions:} U.S. Const. art. I, \S~8, cl. 3 (Dormant Commerce Clause); U.S. Const. amend. XXI, \S~2; Wilson Act, 27 U.S.C. \S~121; Webb-Kenyon Act, 27 U.S.C. \S~122; Tenn. Code Ann. \S~57-3-204(b)(2)--(3)
    \item \textbf{Primary Sources:} \href{https://www.supremecourt.gov/opinions/18pdf/18-96_5i36.pdf}{SupremeCourt.gov Slip Opinion}; \href{https://supreme.justia.com/cases/federal/us/588/18-96/}{Justia}; \href{https://www.law.cornell.edu/supremecourt/text/18-96}{Cornell LII}
\end{itemize}

\section{Facts and Procedural History}

Tennessee regulates alcohol distribution through a three-tiered system separating producers, wholesalers, and retail package stores. Under Tenn. Code Ann. \S~57-3-204(b), Tennessee imposed steep durational-residency requirements on anyone seeking to operate a retail liquor store: (1) an individual applying for an initial license had to be a ``bona fide resident'' of Tennessee for the prior two years; (2) renewing a license after one year required 10 consecutive years of Tennessee residency; and (3) a corporate applicant could not receive a license unless 100 percent of its officers, directors, and stockholders satisfied those residency requirements, effectively barring all publicly traded corporations.

In 2012, the Tennessee Attorney General concluded that these residency mandates violated the Dormant Commerce Clause, prompting the Tennessee Alcoholic Beverage Commission (TABC) to stop enforcing them. After the Tennessee General Assembly added a legislative-intent clause invoking public ``health, safety and welfare,'' the Attorney General reaffirmed his opinion of unconstitutionality, and the TABC maintained its nonenforcement policy.

In 2016, two applicants that did not meet the two-year residency requirement applied for retail liquor licenses: (1) Tennessee Fine Wines and Spirits, LLC, doing business as Total Wine Spirits Beer \& More (Total Wine), a Tennessee LLC owned by Maryland residents seeking to open a store in Knoxville; and (2) Affluere Investments, Inc., doing business as Kimbrough Fine Wine \& Spirits, owned by Doug and Mary Ketchum, who had recently moved from Utah to Memphis to obtain cleaner air and medical care for their daughter (\href{https://ij.org/case/tennessee-wine-and-spirits-retailers-association-v-thomas/}{Institute for Justice}). When TABC staff recommended approving both applications, the Tennessee Wine and Spirits Retailers Association (a trade group of incumbent in-state liquor stores) threatened to sue. The TABC's executive director filed a declaratory judgment action in state court, which was removed to the U.S. District Court for the Middle District of Tennessee. Relying on \href{https://supreme.justia.com/cases/federal/us/544/460/}{\textit{Granholm v. Heald}, 544 U.S. 460 (2005)}, the District Court held the residency rules unconstitutional, and the TABC issued the licenses without appealing. The Association appealed to the Sixth Circuit, where a divided panel (with Judge Sutton dissenting in part) struck down all three residency provisions. The Association petitioned for certiorari solely on the two-year initial residency requirement.

\section{Issue(s) Presented}

Whether \S~2 of the Twenty-first Amendment empowers a State, notwithstanding the Dormant Commerce Clause, to enforce a two-year durational-residency requirement on individuals and corporate officers/directors applying for an initial retail liquor store license.

\section{Holding}

\textbf{No; affirmed} (7--2). Tennessee's two-year durational-residency requirement for retail liquor license applicants facially discriminates against nonresidents in violation of the Dormant Commerce Clause and is not saved by \S~2 of the Twenty-first Amendment. Because \S~2 does not shield protectionist alcohol regulations lacking a predominant, evidence-backed public health or safety justification, the requirement is unconstitutional.

\section{Reasoning}

Justice Alito's opinion for the seven-Justice majority proceeded in four steps: reaffirming the Dormant Commerce Clause, situating \S~2 in pre-Prohibition history, extending \textit{Granholm}'s nondiscrimination rule beyond producers to retail distribution, and applying an evidence-based \S~2 test.

\subsection{Vitality of the Dormant Commerce Clause}
Addressing textualist critiques, Alito emphasized that curbing state economic protectionism was a central catalyst for the 1787 Constitutional Convention and that, given narrow interpretations of the Import-Export and Privileges and Immunities Clauses, the Dormant Commerce Clause remains the primary constitutional bulwark against interstate trade barriers. Outside the alcohol context, Tennessee's two-year residency requirement would plainly fail strict scrutiny because it facially favors Tennesseans over nonresidents.

\subsection{History and Unified Reading of Section 2}
Section 2 prohibits ``[t]he transportation or importation into any State \ldots{} for delivery or use therein of intoxicating liquors, in violation of the laws thereof.'' Alito rejected a hyper-literal reading that would let \S~2 override every earlier constitutional guarantee (such as Equal Protection, Free Speech, or the Ex Post Facto Clause). Instead, \S~2 must be read as part of a unified constitutional scheme that ``constitutionalized'' the pre-Prohibition balance struck by the Wilson Act (1890) and Webb-Kenyon Act (1913). Those statutes eliminated the ``original-package'' loophole so dry States could regulate imported alcohol on equal terms with domestic alcohol under bona fide police powers, but they never authorized purely protectionist discrimination against out-of-state citizens or goods. While early post-Repeal decisions such as \textit{State Bd. of Equalization v. Young's Market Co.}, 299 U.S. 59 (1936), suggested plenary state authority, modern cases---\href{https://supreme.justia.com/cases/federal/us/468/263/}{\textit{Bacchus Imports, Ltd. v. Dias}, 468 U.S. 263 (1984)}, \textit{Healy v. Beer Institute}, 491 U.S. 324 (1989), and \href{https://supreme.justia.com/cases/federal/us/544/460/}{\textit{Granholm v. Heald}, 544 U.S. 460 (2005)}---firmly established that economic protectionism is not a legitimate \S~2 interest.

\subsection{No Producer-Versus-Retailer Carveout}
The Association argued that \textit{Granholm}'s nondiscrimination principle applied only to out-of-state \textit{products and producers}, leaving in-state wholesale and retail distribution immune. Alito rejected this distinction as both textually backwards (since \S~2's text speaks of ``transportation or importation,'' not domestic store licensing) and doctrinally unsound. While \textit{Granholm} described the basic three-tier model as ``unquestionably legitimate,'' it did not immunize every discriminatory add-on a State attaches to that model. Durational-residency rules are not an ``essential feature'' of a three-tier system; many three-tier States operate without them.

\subsection{Application of the Predominant-Effect Test}
Under \S~2, States enjoy greater regulatory leeway over alcohol than over ordinary goods, so the Court asks whether a challenged law ``can be justified as a public health or safety measure or on some other legitimate nonprotectionist ground''---not by ``mere speculation,'' but by ``concrete evidence.'' Where the law's ``predominant effect'' is protectionism, \S~2 affords no shield. Tennessee's two-year requirement failed every asserted rationale: (1) amenability to process is easily secured by requiring a registered agent; (2) fitness screening can be achieved via background checks (indeed, the TABC does not even investigate prospective applicants during their two-year wait); (3) on-site inspections, audits, and license revocation ensure oversight of brick-and-mortar stores; and (4) promoting ``responsible consumption'' through community familiarity was unsupported by evidence and fatally underinclusive (a resident of Memphis, 500 miles away, could open a store in Bristol, Tennessee, while a resident right across the border in Bristol, Virginia, could not, and Tennessee imposed no residency rule on bar owners).

\section{Concurring and Dissenting Opinions}

There were no concurring opinions. Justice Gorsuch, joined by Justice Thomas, dissented.

Gorsuch argued that the Dormant Commerce Clause is an unwritten, ``peculiar'' judicial implication that Congress can override---and that the people themselves overrode for alcohol by constitutionalizing the Webb-Kenyon Act in \S~2 of the Twenty-first Amendment. Recounting over 150 years of state liquor residency laws (dating back to New Hampshire in 1834 and Missouri in 1845) and 18 post-Repeal state residency statutes, Gorsuch contended that \S~2 embodied a federalist compromise allowing each State to restrict competition, raise liquor prices, and tailor distribution to local norms free from judicial free-trade policing.

Gorsuch further argued that \textit{Granholm} intentionally drew a line between out-of-state \textit{products/producers} and in-state \textit{distribution tiers} to preserve traditional residency and three-tier rules. By erasing that line and announcing an amorphous ``predominant effect'' standard, Gorsuch warned, the majority left lower courts with no workable guidance on how much health-and-safety benefit suffices or whether basic in-state physical-presence requirements for retailers can survive.

\section{Key Quotations}

All quotations are verified verbatim from the official slip opinion (\href{https://www.supremecourt.gov/opinions/18pdf/18-96_5i36.pdf}{SupremeCourt.gov}; \href{https://www.law.cornell.edu/supremecourt/text/18-96}{Cornell LII}):

\begin{quote}
``Section 2 was adopted as part of the scheme that ended prohibition on the national level. It gives each State leeway in choosing the alcohol-related public health and safety measures that its citizens find desirable. But \S~2 is not a license to impose all manner of protectionist restrictions on commerce in alcoholic beverages.'' \\
--- Justice Alito, slip op. at 2 (\href{https://www.supremecourt.gov/opinions/18pdf/18-96_5i36.pdf}{SupremeCourt.gov})
\end{quote}

\begin{quote}
``At issue in the present case is not the basic three-tiered model of separating producers, wholesalers, and retailers, but the durational-residency requirement that Tennessee has chosen to impose on new applicants for liquor store licenses. Such a requirement is not an essential feature of a three-tiered scheme.'' \\
--- Justice Alito, slip op. at 28 (\href{https://www.supremecourt.gov/opinions/18pdf/18-96_5i36.pdf}{SupremeCourt.gov})
\end{quote}

\begin{quote}
``Section 2 gives the States regulatory authority that they would not otherwise enjoy, but as we pointed out in \textit{Granholm}, `mere speculation' or `unsupported assertions' are insufficient to sustain a law that would otherwise violate the Commerce Clause. \ldots{} Where the predominant effect of a law is protectionism, not the protection of public health or safety, it is not shielded by \S~2.'' \\
--- Justice Alito, slip op. at 33, 588 U.S. at 539--540 (\href{https://www.supremecourt.gov/opinions/18pdf/18-96_5i36.pdf}{SupremeCourt.gov})
\end{quote}

\begin{quote}
``Like it or not, those who adopted the Twenty-first Amendment took the view that reasonable people can disagree about the costs and benefits of free trade in alcohol. They left us with clear instructions that the free-trade rules this Court has devised for `cabbages and candlesticks' should not be applied to alcohol. \ldots{} Under the terms of the compromise they hammered out, the regulation of alcohol wasn't left to the imagination of a committee of nine sitting in Washington, D. C., but to the judgment of the people themselves and their local elected representatives.'' \\
--- Justice Gorsuch, dissenting, slip op. at 15 (\href{https://www.supremecourt.gov/opinions/18pdf/18-96_5i36.pdf}{SupremeCourt.gov})
\end{quote}

\section{Doctrinal Significance and Modern Legacy (Through 2026)}

\textit{Tennessee Wine} reshaped both Twenty-first Amendment doctrine and broader Dormant Commerce Clause jurisprudence. First, it decisively settled the post-\textit{Granholm} circuit split over whether the Dormant Commerce Clause's antidiscrimination principle reaches the wholesale and retail tiers of state alcohol markets: it does, and States cannot insulate protectionist licensing barriers by invoking \S~2. Second, Justice Alito's robust defense of the Dormant Commerce Clause's historical legitimacy anchored the Court's antidiscrimination core even as a fractured Court later narrowed the extraterritoriality and \textit{Pike} balancing branches in \href{https://supreme.justia.com/cases/federal/us/598/21-468/}{\textit{National Pork Producers Council v. Ross}, 598 U.S. 356 (2023)}.

Through 2026, however, Justice Gorsuch's dissent proved prescient: lower courts have deeply fractured over how \textit{Tennessee Wine} applies to state laws allowing \textit{in-state} wine retailers to ship directly to consumers while barring \textit{out-of-state} retailers unless they maintain an in-state physical storefront. Two rival methodologies emerged across the circuits:
\begin{itemize}
    \item \textbf{The ``Essential Feature'' Shortcut:} The Fourth, Eighth, and Ninth Circuits---most prominently in \href{https://law.justia.com/cases/federal/appellate-courts/ca9/23-16148/23-16148-2025-03-04.html}{\textit{Day v. Henry}, 152 F.4th 961 (9th Cir. 2025)} (amended opinion, Sept. 2025)---upheld in-state physical-premises requirements for wine retailers by treating physical presence as an ``essential feature'' of the three-tier system, largely bypassing strict scrutiny of nondiscriminatory alternatives.
    \item \textbf{The ``Concrete Evidence'' Inquiry:} The Sixth Circuit in \textit{Block v. Canepa}, 74 F.4th 400 (6th Cir. 2023), and subsequent remand proceedings insisted that courts must weigh concrete record evidence under \textit{Tennessee Wine}'s predominant-effect test rather than deferring to a State's bare invocation of three-tier architecture (\href{https://libationlawblog.com/2026/06/24/ohio-wine-shipping-supreme-court-cert-petition/}{Libation Law Blog, June 2026}).
\end{itemize}
On May 18, 2026, the Supreme Court denied certiorari without comment in both \textit{Day v. Henry} (No. 25-788) and \textit{Chicago Wine Co. v. Braun} (No. 25-844) (\href{https://libationlawblog.com/2026/05/20/supreme-court-denies-cert-day-henry-chicago-wine-braun-retailer-shipping/}{Libation Law Blog, May 20, 2026}), leaving the circuit split unresolved as renewed Ohio retailer-shipping petitions reached the Court in late 2026.

\section{Pedagogical Critiques and Exam Hypotheticals}

\subsection{Critiques}
\begin{itemize}
    \item \textbf{Doctrinal Indeterminacy:} \textit{Tennessee Wine} creates a hybrid standard---looser than strict Dormant Commerce Clause scrutiny (allowing ``leeway'' for alcohol) yet far stricter than rational-basis review (demanding ``concrete evidence'' that health or safety predominates over protectionism)---without specifying how courts should weigh mixed-motive laws.
    \item \textbf{Tension with Public-Health Temperance Theory:} As Gorsuch noted, limiting market entry and raising retail liquor prices can reduce alcohol consumption; Alito's opinion never cleanly explains when higher prices count as legitimate temperance regulation versus illegitimate economic protectionism.
    \item \textbf{Textualism Versus Constitutional Structure:} The 7--2 split illustrates a classic methodological divide: Gorsuch and Thomas prioritized the enacted text of \S~2 and its 1933 original public meaning over an implied Commerce Clause doctrine, whereas Alito's majority harmonized \S~2 with the Constitution's broader anti-Balkanization structure.
\end{itemize}

\subsection{Exam Hypothetical and ``Opposite Side'' Contrast}
\textbf{Hypothetical (Discriminatory Retail Licensing):} State X enacts a statute allowing only individuals who have owned real property in State X for three years to hold a craft-beer wholesale or retail license, asserting that property owners are ``more invested in preventing underage drinking.'' Out-of-state applicant Z challenges the law. \textit{Analysis:} Under \textit{Tennessee Wine}, the statute is \textbf{unconstitutional}. It facially discriminates against nonresidents, is not an essential feature of a three-tier system, and rests on ``mere speculation'' rather than concrete evidence, while nondiscriminatory alternatives (age-verification audits, clerk training, license revocation) readily advance public safety.

\textbf{``Opposite Side'' Contrast:} Now suppose State X eliminates all residency and property-ownership requirements and instead enacts a \textit{facially neutral} statute capping the total number of retail liquor licenses per county at one per 10,000 residents, imposing a uniform minimum retail price per ounce of alcohol, and requiring every licensed retailer (in-state or out-of-state owned) to operate a brick-and-mortar storefront in State X subject to unannounced state health-and-tax inspections, backed by empirical studies tying outlet density and price floors to reduced alcohol-related hospitalizations. Here, State X holds a far stronger defense: the density cap and uniform price floor are nondiscriminatory public-health measures squarely within \S~2's ``leeway'' (which Alito expressly endorsed, slip op. at 36), and under post-\textit{Tennessee Wine} circuit precedent such as \href{https://law.justia.com/cases/federal/appellate-courts/ca9/23-16148/23-16148-2025-03-04.html}{\textit{Day v. Henry} (9th Cir. 2025)}, an evenhanded physical-presence requirement tied to three-tier inspection and tax enforcement is far more likely to survive \S~2 scrutiny than a durational-residency bar.

\end{document}
